\section{Alignments in sparse prime steps}\label{sec:sparse}

We now estimate the alignment terms that arise during a sparse-prime
update in an existing component.  The estimates use only the previously
constructed tables, their core equivariance, and the static birth
prescriptions.  In particular, they do not require a choice of the next
operational partition.  Besides paying for alignments outside a very small
distance scale, we construct deterministic interval trains covering the
segments by which accounting classes will be joined.  These trains are
the input to the partition construction in Section~\ref{sec:partitions}.

Fix a component at height $Y$ and a common prime $q$ in its current bin.
For two centres $L=L_i$ and $M=L_j$, retain the notation
\[
 G=(L,M),\qquad U=L/G,\qquad T=M/G,\qquad
 \Theta_{ij}=G r_{\min},\qquad K=[L,M]r_{\max}.
\]
At a successful plus endpoint the denominator has the form
\[
 V'=\frac{L}{dq},\qquad\text{or}\qquad W'=\frac{M}{eq},
\]
where $d,e$ are processed deficits below $q$.  Write $a,b$ for the
remaining masks above $q$, and put
\begin{equation}\label{sp:signatures}
 f_0=(Td,Ue),\qquad s_0=Td/f_0,\qquad t_0=Ue/f_0.
\end{equation}
Thus the successful positions are $(z+a\gamma)/V'$ and
$(z'+b\gamma)/W'$.  Throughout this section their mask coefficients are
the \emph{product} of the two after-$q$ probabilities, including when the
two masks happen to be equal.  They are not the single mask probability
used in clipping.

The endpoint train test below has the following concrete form.  For each
successful denominator $V'$, a deterministic list contains at most
$Z_Y=\exp(\exp(h_aY))$ offsets $y_0\pmod1$.  An endpoint of width $r$
passes if it lies within $r\exp(-.4\exp(\sigma_gY))$ of some point
$(z+y_0)/V'$ with $z\in\mathbb Z$ and a listed offset.  The lists depend
only on the centre, prime and processed mask, not on the actual table
history.  Construction~\ref{sp:danger-trains} specifies them using
rational residues and quantile marks.

\begin{proposition}[Sparse-step estimates]\label{sp:sparse-estimates}
Consider any finite table history up to a sparse-prime update that satisfies
the centre and good-mask prescriptions and the preceding chronological
core-equivariance rules.  Suppose its operational cells have the span and
padding required in Section~\ref{sec:tables}.  Uniformly in that history
and in the slot cap, the following costs, for a fixed centre pair and bin,
are at most
\[
 C\Theta_{ij}e^{-cY}\tau(UT)^C,
\]
where $c>0$ and $C$ are fixed independently of the block:
\begin{enumerate}
 \item the width-weighted product-mask counts of actual successful plus
 pairs at distances between
 $r_{\max}\exp(-\exp(\sigma_gY))$ and $C_0r_{\max}$;
 \item the corresponding counts of before-$q$ plus pairs at distance at
 most $C_0r_{\max}$, multiplied by $1/q^2$, with after-$q$ mask coefficients;
 \item for equal widths $r$, the actual successful plus pairs at distance
 at most $r\exp(-\tfrac12\exp(\sigma_gY))$ that fail the deterministic
 train test of Construction~\ref{sp:danger-trains} at either endpoint.
\end{enumerate}
Here each count is summed over the primes of the bin and uses the factor
$r_{\min}$, or $r$ in the third case.  The first and third counts may use
the actual raw arrival weights immediately after the update; the second
uses the before-update raw plus weights.  The constant $C_0$ can be any
fixed positive constant.

Join accounting classes by all actual surviving equal-width successful
plus pairs at the relaxed scale in the third case that pass both train
tests.  Each connecting segment is covered by the deterministic train
history, which is invariant under the post-generation core translations.
Summing the displayed errors over centre pairs, bins, and components is
bounded.  Consequently this construction proves
Proposition~\ref{tab:sparse-input}.
\end{proposition}

We prove the proposition in stages.  All upper estimates below may drop
assignment restrictions and adaptive filters when indicated.  The actual
arrival weights are bounded by the static factors retained in those upper
estimates, since a surviving label follows same-location ancestors with
nonincreasing raw weight.

\subsection{Bare counts and switching}

A pair at distance at most $C_0r_{\max}$ satisfies
\begin{equation}\label{sp:determinant}
 \left|s_0(z+a\gamma)-t_0(z'+b\gamma)\right|
 \le \frac{C_0K}{qf_0}.
\end{equation}
By \eqref{eq:component-path}, the right side is less than $1/2$ after
increasing the fixed lower cutoff.  Thus there is at most one possible
integer $s_0z-t_0z'$ for fixed masks.  The gcd of the two successful
denominators is $Gf_0/(qde)$, so their width-weighted bare count is at most
\begin{equation}\label{sp:bare-count}
 \Theta_{ij}\frac{f_0}{qde}.
\end{equation}
The same primewise gcd-ratio calculation used in
Lemma~\ref{dir:ultra-switch} bounds the unrestricted sum of $f_0/(de)$
by $\exp(CY)\tau(UT)^C$.  Outer masks still integrate as independent
probabilities.

For the second assertion of Proposition~\ref{sp:sparse-estimates}, use the
before-$q$ denominators instead.  The number of determinant integers is
$O(1+K/f_0)$, and the width-weighted pre-count is therefore at most
\[
 C\Theta_{ij}\frac{K+f_0}{de}.
\]
Multiply by $q^{-2}$ and sum over $q\ge\exp(\exp Y)$.  The resulting
double-exponential saving absorbs both the crude mask cost and the bound
on $K$ from \eqref{eq:component-path}.  This proves the required
$\Theta_{ij}e^{-cY}\tau(UT)^C$ estimate.

For the remaining assertions, use the switching argument with budget
$u=D''Y$, where $D''$ is sufficiently large absolutely.  Test divisors
$h\mid(d,s_0)$ in the ordinary dyads between
$\exp(C_1u)$ and $\exp(3C_1u)$, with the cutoff rounded as in
Lemma~\ref{dir:ultra-switch}, and test every such prime in dyads above the
lower cutoff.  Make the symmetric tests on $(e,t_0)$.

On a fiber with $d/h,e,a,b,q$ fixed, stripping $h$ leaves $f_0$ unchanged.
Indeed, at a prime $l$, write $A_l=\nu_l(Td)$, $B_l=\nu_l(Ue)$,
and $j_l=\nu_l(h)$.  The divisibility $h\mid(d,s_0)$ gives
$j_l\le A_l-\min(A_l,B_l)$, so
\[
 \min(A_l-j_l,B_l)=\min(A_l,B_l).
\]
It replaces $s_0$ by $s_0/h$ and extracts a factor
$1/h$ from \eqref{sp:bare-count}.  The tolerance in
\eqref{sp:determinant} stays fixed, while $h$ occurs linearly in the
phase.  The integer rotation estimate therefore applies on this fiber.
The stripped row need not remain assigned or good, because those
restrictions have been discarded in this nonnegative upper bound.
Summing its exceptional alternatives over the tested dyads, stripped
masks, and current primes gives
\begin{equation}\label{eq:sparse-switch}
 \ll \Theta_{ij}e^{-cD''Y}\tau(UT)^C.
\end{equation}
There are only $\exp(O(Y))$ prime-size dyads below the end of the bin;
the divisor-test dyads are fewer.  The crude mask and reciprocal-prime
sums have the same $\exp(O(Y))$ allowance.  These costs are absorbed by
the choice of $D''$, leaving the divisor factor displayed in
\eqref{eq:sparse-switch}.

The argument is also valid if the tolerance in
\eqref{sp:determinant} is multiplied by one of the smaller gap factors
used below.  Outside the union of these failures, any available switch
gives coupled, unique rational approximants $A_i,A_j$ to
$s_0a\gamma,t_0b\gamma$.  Their denominators are at most $e^{O(u)}$,
their errors are at most $e^{O(u)}$ times the tolerance, and
\[
 A_i-A_j=-(s_0z-t_0z').
\]
Division by a switched divisor retains a denominator bounded by
$e^{O(u)}$.  In particular a sufficiently large switched prime divides
the relevant numerator.  The divisor-assembly argument in
Lemma~\ref{dir:ultra-switch}, together with the good-mask prime-power
bound, then gives the corresponding small denominators after division
by the full signatures.  If no switch is available, the signatures
themselves are at most $\exp(\exp(h_*Y))$.  Indeed
$s_0/(d,s_0)\mid T$ and $t_0/(e,t_0)\mid U$, the unmatched prime-power
excess is bounded by the purge, and a larger uncancelled squarefree
product would contain a divisor in the tested range.  The remaining
centre-ratio cost is covered by \eqref{eq:component-path}.

\subsection{Neutralizing earlier dense bins}

Call the intermediate-distance range in the first assertion of
Proposition~\ref{sp:sparse-estimates} the \emph{off-gap} range.
The crude factor $\exp(CY)$ is too large for its sum.  We now
remove the contribution of any earlier dense bins near the current
height by using the birth prescriptions.

The centre exponents agree in every bin of height $Y'\ge\sigma Y$,
by \eqref{eq:component-path}.  If no such bin is dense, the reciprocal
sum through the current bin is already $O((\sigma+\kappa)Y)$: bins
below $\sigma Y$ have the absolute harmonic bound, and all remaining
bins are sparse.  Otherwise the last dense height $Y_d\ge\sigma Y$
is common to both centres, since bins after the start are sparse.
The starts are also common.  On medium type,
\[
 \frac{\sigma Y}{D_0}\le B_i,B_j\le\frac{Y}{A_0};
\]
the retained band separation forces $B_i=B_j$.  On high type each start
is $2Y_d$.  On medium type all bins with
$\sigma Y\le Y'<Y_0$ are ordinary tag bins.

Write $d=d^0u_i$ and $e=e^0u_j$, where $d^0,e^0$ are below this common
start.  At birth, the ancestors of the current successful points had
phase and numerator
\[
 \alpha_i^0=u_iqa\gamma,\qquad z_i^0=u_iqz,
 \qquad
 \alpha_j^0=u_jqb\gamma,\qquad z_j^0=u_jqz'.
\]
Consequently the birth integer on the first side is
$\mathcal L_i=\lfloor dqa\gamma\rfloor$, and similarly on the second.
Switch every tag mismatch prime in the indicated bins of height at
least $\sigma Y$; on high type use its tag bin $Y_d$.  Such a prime is
eligible for \eqref{eq:sparse-switch}: centre exponents agree, its good
deficit has exponent at most one, and it divides the relevant signature
when it occurs on just one side.

If a mismatch occurs, form the predicted tag lifts
\begin{equation}\label{sp:predicted-tags}
 F_i^{\mathrm{pred}}=\mathcal L_i-\frac{qf_0}{T}A_i,
 \qquad
 F_j^{\mathrm{pred}}=\mathcal L_j-\frac{qf_0}{U}A_j.
\end{equation}
Their log sizes and log denominators are at most $\exp(h_*Y)$.
For the size bound, the large phase cancels:
$(qf_0/T)s_0a\gamma=dqa\gamma$, so the remaining size is controlled
by the approximation error times $qf_0/T$, plus the floor error.
The denominator bound follows from the same approximant bounds and
\eqref{eq:component-path}.  These bounds fit the recognized tag heights
because $h_*\ll\sigma h_H$.

At every included anchor on the first side,
$F_i^{\mathrm{pred}}\equiv\mathcal L_i$.  A common omission supplies
the prime through $f_0$; a mismatched omission supplies it through the
approximant numerator.  The denominators are units there.  Thus every
included anchor reconstructs this prediction.  The coupled exact
identity gives
\[
 T(d^0z_i^0+\mathcal L_i-F_i^{\mathrm{pred}})
 =U(e^0z_j^0+\mathcal L_j-F_j^{\mathrm{pred}}).
\]
Reduction at a mismatch prime contradicts the full side's tag exclusion
if both tag masks in that bin are nonempty.  Hence each such bin has
identical masks or one empty mask.  On high type the occupancy purge
forces the tag masks to be the same nontrivial mask $m$.

\paragraph{Earlier low bins on high type.}
On high type it remains to handle bins
$\sigma Y\le Y'<Y_d$.  Use the birth signatures
\[
 f_B=(Td^0,Ue^0),\qquad s_B=Td^0/f_B,
 \qquad t_B=Ue^0/f_B.
\]
Their birth determinant has tolerance
$C_0K/f_B\le C_0K/m$, which is sufficiently small by the size of a
prime in the common nontrivial tag mask and
\eqref{eq:component-path}.  Make the same divisor and prime switches
on $(d^0,s_B)$ and $(e^0,t_B)$, holding the birth phases fixed.
Stripping such a factor preserves both $f_B$ and $f_0$, since its primes
are below the common start: at each of them $u_i,u_j$ have zero
valuation, so the preceding primewise minimum calculation applies
to both gcds.  Fixing $q,u_i,u_j,a,b,e^0,d^0/h$ fixes both birth
phases during the rotation test.  The exact rescaling is
\[
 s_Bu_iq=\frac{qf_0}{f_B}s_0,
 \qquad t_Bu_jq=\frac{qf_0}{f_B}t_0.
\]
Thus the birth determinant is the current determinant multiplied by
$qf_0/f_B$, and its tolerance is $C_0K/f_B$, unchanged on the
stripping fiber.  The charge is still the current success count
\[
 \Theta_{ij}\frac{f_0}{q d^0u_i e^0u_j}:
\]
stripping extracts $1/h$ without removing its factor $1/q$.
The rotation error is uniform in the fixed birth phases, even though
those phases vary when $q$ is subsequently summed.  Hence
\eqref{eq:sparse-switch}, including its $1/q$ summation, still applies.

When a switch occurs, divide the resulting coupled approximants by
$s_B,t_B$.  Their denominators are at most
$\exp(\exp(h_*Y))$: use
$s_B/(s_B,d^0)\mid T$, the analogous divisibility on the other side,
the prime-power purge, and the divisor-assembly consequences of the
switches.  If $\psi_i^0=(L/d^0)r_i$ is the first birth numerator
width, then
\[
 \frac{C_0K/f_B}{s_B}
 =\frac{C_0Lr_{\max}}{d^0}
 =C_0\psi_i^0\frac{r_{\max}}{r_i}.
\]
Thus its approximation error divided by $\psi_i^0$ is at most
$e^{O(D''Y)}r_{\max}/r_i\le\exp(\exp(h_*Y))$, and similarly
on the other side.  The models concern the transported birth phases
$\alpha_i^0=u_iqa\gamma$, not the current phases $a\gamma$;
no division of a birth-model denominator by $u_iq$ is required.
Since $Y_d\ge\sigma Y$ and $h_*\ll\sigma h_H$, the high birth
models therefore qualify
and yield exact equality for the birth determinant.  At a mismatch
prime in these bins one signature is divisible by the prime and the
model denominators are units.  Exact equality then violates the hard
exclusion on the full side.  This prime lies above the hard threshold
by $h_*\ll\sigma h_H$ and $h_H\ll\sigma$.
If there is no switch, the signatures themselves are at most
$\exp(\exp(h_*Y))$, so they cannot contain a mismatch prime of height
at least $\sigma Y$.  Thus those masks are identical as well.

\paragraph{Local factors after reconstruction.}
We can now sum the neutralized coordinates using only static upper
weights, dropping other processed filters.  On an identical bin, the
gcd-ratio coordinate is $1/p$ at a common omission and $1$ at a full
factor.  If a common residual power remains, retain one side's scalar
or hard factor.  Its mean is $R_p$ on the residual common digit, since
$s_0,t_0,u_iq$ are units at $p$.  The summed local cost is at most
\[
 R_p+\frac1p R_p^{[\text{common exponent at least }2]}\le1.
\]
For a low static prime on high type, retain any one active hard
exclusion, or the actual scalar if no exclusion is active; discard the
other low tests for this upper count.

In a one-empty tag bin the empty side uses the ordinary scalar factors
$R_p$ prescribed at birth.
Neither-omitted coordinates cost at most $R_p^2$ on averaging the
common digit.  The two possible one-sided omissions have bare cost
$1/p$ each.  Even allowing both choices bounds the total by
$1+O(p^{-2})$.  Fix masks and anchor choices before using digit CRT,
as in Lemma~\ref{dir:tag-fold}.  There are only a bounded number of
bins between $\sigma Y$ and $Y$, so summing their identical/empty
arrangements costs a constant.  All unneutralized bins together now
cost only
\begin{equation}\label{sp:reduced-harmonic-cost}
 \exp\bigl(C(\sigma+\kappa)Y\bigr)\tau(UT)^C,
\end{equation}
with absolute coefficients in the harmonic exponent.

The earlier dense bins have therefore been removed from the harmonic
cost.  We next sum the remaining factor $1/q$ over the distance window
for off-gap pairs, then estimate failures of the endpoint train tests
for the tightly aligned equal-width pairs.

\subsection{The off-gap prime window}

We retain the factor $1/q$ in \eqref{sp:bare-count}.  To make the
remaining variables independent of the position of $q$ inside its bin,
extend $d,e$ through the end of that bin.  Write $a=a_Ya_>$ and
$b=b_Yb_>$ for their within-bin and above-bin parts.  Let $I_{>q}$
be the common prime-power part above $q$, and let $I_>$ be its part
above the bin.  The actual outer laws
satisfy the pointwise bound
\[
 \pi_{I_{>q}}(a)\pi_{I_{>q}}(b)
 \le \frac{\pi_{I_>}(a_>)\pi_{I_>}(b_>)}{a_Yb_Y},
\]
because $\pi_{p^\nu}(p^j)=R(p^{\nu-j})p^{-j}\le p^{-j}$.
Pad $a_Y,b_Y$ to arbitrary divisors of the whole bin product.  At
each common bin prime $p^\nu$, the total cost of the four extended
coordinates is bounded by
\[
 \left(\sum_{j,k=0}^{\nu}p^{-\max(j,k)}\right)
 \left(\sum_{j=0}^{\nu}p^{-j}\right)^2
 \le 1+\frac Cp.
\]
Thus padding costs at most $\exp(C\kappa Y)$.  Actual allowed
choices still require that $q$ be absent from all four masks and
that their before/after supports have the proper order; dropping
these restrictions later only enlarges the prime window.

We also need a bound for each fixed padded tuple that is independent
of $q$, not merely a bound after summing masks at a fixed $q$.
At a neutralized earlier prime, write $j,k\in\{0,1\}$ for its two
deficits.  After static digit averaging, the identical-bin mode is
bounded by the local majorant
\[
 Q_p(0,0)=R_p,\qquad
 Q_p(1,1)=p^{-1}R_p^{[\nu\ge2]},\qquad
 Q_p(1,0)=Q_p(0,1)=0.
\]
In the one-empty mode use instead
\[
 Q_p(0,0)=R_p^2,\qquad
 Q_p(1,0)=Q_p(0,1)=p^{-1},\qquad Q_p(1,1)=0.
\]
The latter permits either one-sided omission separately at each
prime and hence enlarges a bin with one globally empty mask.
The coordinate sums are at most $1$ and $1+O(p^{-2})$, respectively.
These majorants hold for every fixed anchor choice: their derivation
uses only one test's mean, or an ordinary scalar $R_p$ together with the other
test's mean, and is independent of the forbidden offsets.  Thus
the birth offsets may vary with $q$ without changing $Q_p$.
Adaptive weights have already been dropped.  Sum over the bounded
number of bin-mode patterns, not over a separate union of modes at
each prime.  Together with the unneutralized gcd coordinates and
the padded outer-law bound, this gives a $q$-independent majorant
for every fixed tuple, whose sum is bounded by
\eqref{sp:reduced-harmonic-cost} with the padding allowance.

For a fixed choice of these integers, the normalized distance of an
actual cohit is
\[
 \frac{\text{distance}}{r_{\max}}
 =\|s_0a\gamma-t_0b\gamma\|\frac{qf_0}{K}.
\]
The nearest determinant is unique.  If the norm vanishes the pair is
not an off-gap pair; otherwise the required distance range restricts
$q$ to a multiplicative window whose logarithmic width is
$O(1+\exp(\sigma_gY))$.  Elementary prime counting in ordinary dyads
therefore gives
\[
 \sum_{q\text{ in this window and bin}}\frac1q
 \ll \exp\bigl(-(1-\sigma_g)Y\bigr),
\]
since $\log q\ge\exp Y$.  The neutralized digit factors above are
independent of this $q$ summation: they belong to earlier bins.  We may
first enlarge the actual allowed prime support to this whole window
and then sum those local bounds.  Combining with
\eqref{sp:reduced-harmonic-cost}, and using the choices of $\sigma$ and
$\kappa$, proves the first estimate of
Proposition~\ref{sp:sparse-estimates}.

\subsection{Train coverage of tightly aligned pairs}

Only the relaxed tight equal-width scale remains.  We construct its
exceptional set by means of lists fixed before the table history is
used.  The purpose of the construction is twofold: its failures have
small product-mask cost, while every accepted connecting segment has
a deterministic spatial cover that can be protected at future joins.

\begin{construction}[Deterministic danger trains]\label{sp:danger-trains}
For a centre $i$, current prime $q$ in bin $Y$, and processed deficit
$d$ such that $dq$ is the processed portion of some good mask, put
$V'=L_i/(dq)$ and use a list of residues $y_0\pmod1$ of size at most
\[
 Z_Y=\exp(\exp(h_aY)).
\]
Include in the list every rational residue of denominator at most
$\exp(\exp(h_*Y))$.  Also include $0$ and generalized inverse
cumulative quantiles for the distribution of $a\gamma\pmod1$ under
the unrestricted deficit law on the prime powers of $L_i$ above $q$,
using
\[
 N=\left\lceil\exp(\exp(2h_*Y))\right\rceil
\]
equal probability increments.  Repeated quantiles are included only
once.  Each open arc disjoint from the marks then has probability
$O(1/N)$.

The endpoint test requires $x$ to lie within
\[
 r_i\exp\bigl(-.4\exp(\sigma_gY)\bigr)
\]
of a point $(z+y_0)/V'$, for an integer $z$ and a listed residue.
The stored train history uses the same grids with the larger radius
\[
 r_i\exp\bigl(-.3\exp(\sigma_gY)\bigr).
\]
The supply consists of all these grids for the indicated good processed
masks, whether or not any actual arrival uses them.
\end{construction}

The list-size bound follows because the rational residues number at most
$O(\exp(2\exp(h_*Y)))$ and the quantiles number $O(N+1)$, while
$3h_*<h_a$.  The outer law used for the quantiles depends only on the
centre and the current prime, not on matching decisions, partitions, or
which rows survive.  Each train is invariant under the post-$q$ core:
that core divides $V'$, by the construction underlying
\eqref{eq:core-loss}.  Thus this invariance persists at every later
component change.

Consider an actual successful pair of common width $r$ with distance at
most $r\exp(-\tfrac12\exp(\sigma_gY))$.  Apply the switching estimate
\eqref{eq:sparse-switch} at this smaller tolerance.  If a switch is
available outside its paid failures, the recovered rational models for
$a\gamma,b\gamma$ have denominators at most
$\exp(\exp(h_*Y))$.  Their errors are respectively at most
\[
 V'r\exp(-.4\exp(\sigma_gY)),\qquad
 W'r\exp(-.4\exp(\sigma_gY)).
\]
Indeed division by the signature converts the determinant tolerance
to the corresponding phase radius, and all switch losses are absorbed
by the gap between the exponents $-.5$ and $-.4$.  These rationals are
on the lists, so both endpoint tests hold.

If no switch is available, $s_0,t_0\le\exp(\exp(h_*Y))$.
Fix $b$.  The determinant condition restricts $a\gamma\pmod1$ to
at most $s_0$ arcs, each of radius
\[
 O\bigl(V'r\exp(-\tfrac12\exp(\sigma_gY))\bigr).
\]
To see the factor $V'$, divide the determinant tolerance by $s_0$;
the lcm of the successful denominators equals $s_0V'$.
Every such arc has diameter smaller than the phase padding in the
endpoint test.  This also handles marked endpoints: a closed
admissible arc meeting a mark places all its phases within that
diameter of the mark.  Arcs crossing $0$ may be split at $0$, which
is itself marked.  For the possibly atomic phase law, every component
of the complement of the quantile marks has mass at most $1/N$;
otherwise its cumulative distribution would cross some $k/N$
inside that component, placing a generalized inverse mark there.
In particular every atom of mass greater than $1/N$ is marked.
Thus an admissible arc failing the endpoint test is disjoint from
the marks and has unrestricted mass $O(1/N)$.  Failure on the first
side therefore has probability
$O(s_0/N)$, and failure on the second side has probability $O(t_0/N)$.
We may drop assignment restrictions for this estimate, keeping the
two outer laws independent.  These superexponentially small
probabilities absorb the crude mask costs and the current-prime sum.
Together with \eqref{eq:sparse-switch}, they prove the third estimate
of Proposition~\ref{sp:sparse-estimates}.

Declare precisely these pairs failing at least one endpoint test to
be the exceptions to coarsening.  Join every actual surviving
nonexceptional successful pair at the relaxed scale.  Its connecting
segment lies in the stored trains: the distance to an endpoint train
centre is at most the test radius plus the pair distance, which is
smaller than the stored radius for large $Y$.  This remains a valid
historical cover if an intermediate label later disappears.

Finally, every error obtained above is bounded by
$C\Theta_{ij}e^{-cY}\tau(UT)^C$.  Since
$Y\ge A_0\max(B_i,B_j)$, the centre-kernel summation using
\eqref{eq:centre-kernel} and \eqref{eq:centre-lower} applies, retaining
exponential decay for shells and heights.  This proves all assertions
of Proposition~\ref{sp:sparse-estimates}.  We have both the required
within-component energy bounds and a history-independent family of
train intervals covering every class-joining segment; the next section
uses exactly that family to preserve classes when components change.
